\section{总结与延伸}

林俊旸这次演讲展现了 Qwen 团队在 2025 年的全景图，可以从三个维度总结：

\textbf{1. 模型能力的全面进化}
\begin{itemize}
    \item 文本：Reasoning 自然化、119 种语言、超长上下文
    \item 视觉：VL 不降质、GUI 操控、视觉 Reasoning
    \item 生成：图像生成接近真实、编辑精度大幅提升
    \item 语音：Thinker-Talker 端到端架构
    \item 架构：Qwen3-Next Hybrid Linear Attention
\end{itemize}

\textbf{2. 从 Model 到 Agent 的范式转变}
\begin{itemize}
    \item Coding Agent 通过环境交互式 RL 训练
    \item SWE-Bench 67--69 分，能力集成到超大统一模型
    \item 目标从解题转向真实世界的软件工程任务
\end{itemize}

\textbf{3. 开源生态的正向循环}
\begin{itemize}
    \item 小模型（1.8B）源于``让硕士生能毕业''的初心
    \item 多语言源于韩国研究者和巴基斯坦用户的反馈
    \item 图像编辑精度源于开源社区发现的偏移问题
    \item ``如果不是开源社区告诉我们，可能这辈子都想不到''
\end{itemize}

\subsection{拓展阅读}

\begin{itemize}
    \item Qwen 官方博客：\url{https://qwen.ai}
    \item Qwen Chat 体验：\url{https://chat.qwen.ai}
    \item Qwen3-Next 架构博客（含 Attention Gate NeurIPS Best Paper）
    \item SWE-Bench 基准：\url{https://swebench.com}
    \item OpenHands（原 OpenDevin）：开源 Agent 框架
\end{itemize}

%% --- Fin del contenido --- %%

\end{document}
